\section{2012 Applied \& Computational Mathematics}

\subsection{Individual}

\begin{exercise}
\label{applied:2012:1}
In the numerical integration formula
\[
\int_{-1}^{1}f(x)\,dx\approx af(-1)+bf(c),
\]
if the constants $a$, $b$, $c$ can be chosen arbitrarily, what is the highest degree $k$ such that the formula is exact for all polynomials of degree up to $k$? Find the constants $a$, $b$, $c$ for which the formula is exact for all polynomials of degree up to this $k$.
\end{exercise}

\begin{solution}
\textbf{Intuition:} This is a variation of Gaussian quadrature where one node is fixed (at $-1$). We have 3 degrees of freedom ($a, b, c$), so we expect to be able to satisfy 3 equations, corresponding to polynomials $1, x, x^2$, which would mean $k=2$.

\textbf{Step-by-step:}
To be exact for degree $k$, the formula must work for $f(x) = x^j$ for $j=0, \dots, k$.
\begin{enumerate}
    \item $j=0$: $\int_{-1}^1 1 \, dx = 2$. Formula: $a(1) + b(1) = 2 \implies a+b=2$.
    \item $j=1$: $\int_{-1}^1 x \, dx = 0$. Formula: $a(-1) + b(c) = 0 \implies bc = a$.
    \item $j=2$: $\int_{-1}^1 x^2 \, dx = \frac{2}{3}$. Formula: $a(-1)^2 + b(c^2) = \frac{2}{3} \implies a + bc^2 = \frac{2}{3}$.
\end{enumerate}
Substitute $a = bc$ into (1): $bc + b = 2 \implies b(1+c) = 2$.
Substitute $a = bc$ into (3): $bc + bc^2 = \frac{2}{3} \implies bc(1+c) = \frac{2}{3}$.
Using $b(1+c)=2$ in the second equation gives $2c = \frac{2}{3} \implies c = \frac{1}{3}$.
Then $b(1 + 1/3) = 2 \implies b(4/3) = 2 \implies b = \frac{3}{2}$.
Then $a = bc = \frac{3}{2} \cdot \frac{1}{3} = \frac{1}{2}$.

Check for $j=3$: $\int_{-1}^1 x^3 \, dx = 0$.
Formula: $\frac{1}{2}(-1)^3 + \frac{3}{2}(1/3)^3 = -\frac{1}{2} + \frac{3}{2} \cdot \frac{1}{27} = -\frac{1}{2} + \frac{1}{18} \neq 0$.
Thus $k=2$. The constants are $a = 1/2, b = 3/2, c = 1/3$.
\end{solution}

\begin{exercise}
\label{applied:2012:2}
Consider the moving least square approximation of $f(x)$ near $\bar{x}$ given $K$ points $x_k$:
\[
\min_{P_{\bar{x}}\in\Pi_{m}}\sum_{k=1}^{K}|P_{\bar{x}}(x_{k})-f_{k}|^{2}
\]
where $P_{\bar{x}}(x)=\sum_{j=0}^m c_j (x-\bar{x})^j$.

\begin{enumerate}
\item[(a)] Prove there is a unique solution.
\item[(b)] Prove $|c_{i}-\frac{1}{i!}f^{(i)}(\bar{x})|=O(h^{m+1-i})$.
\item[(c)] If points are symmetric, show $O(h^{m+2-i})$ for $i \equiv m \pmod{2}$.
\end{enumerate}
\end{exercise}

\begin{solution}
\textbf{Intuition:} MLS approximates a function locally by a polynomial. Part (a) is standard linear least squares. Part (b) relates the coefficients to Taylor coefficients; the error is dominated by the first neglected term in the Taylor expansion ($x^{m+1}$). Part (c) uses symmetry: for symmetric nodes, the error from the $x^{m+1}$ term vanishes for certain coefficients due to cancellation.

\begin{enumerate}
\item[(a)] Let $V$ be the $K \times (m+1)$ matrix with $V_{kj} = (x_k - \bar{x})^j$. The problem is $\min \|V\mathbf{c} - \mathbf{f}\|^2$. $V$ is a Vandermonde-like matrix. Since $x_k$ are distinct and $K > m$, $V$ has full column rank $m+1$. Thus $V^T V$ is invertible, ensuring a unique solution $\mathbf{c} = (V^T V)^{-1} V^T \mathbf{f}$.

\item[(b)] By Taylor: $f(x_k) = \sum_{j=0}^m \frac{f^{(j)}(\bar{x})}{j!}(x_k - \bar{x})^j + \frac{f^{(m+1)}(\xi_k)}{(m+1)!}(x_k - \bar{x})^{m+1}$.
In vector form: $\mathbf{f} = V \mathbf{c}^* + \mathbf{r}$, where $c^*_j = f^{(j)}(\bar{x})/j!$ and $r_k = O(h^{m+1})$.
The solution is $\mathbf{c} = \mathbf{c}^* + (V^T V)^{-1} V^T \mathbf{r}$.
Let $D = \text{diag}(1, h, \dots, h^m)$. Then $V = \tilde{V} D$ where entries of $\tilde{V}$ are $O(1)$.
$\mathbf{c} - \mathbf{c}^* = D^{-1} (\tilde{V}^T \tilde{V})^{-1} \tilde{V}^T \mathbf{r}$.
Since $\mathbf{r} = O(h^{m+1})$, and $(\tilde{V}^T \tilde{V})^{-1} \tilde{V}^T$ is $O(1)$, we have $c_i - c_i^* = h^{-i} O(h^{m+1}) = O(h^{m+1-i})$.

\item[(c)] For symmetric points, $\sum (x_k - \bar{x})^p = 0$ for odd $p$. In the normal equations, the coupling between even and odd powers vanishes if weights are also symmetric. The error from the $h^{m+1}$ term involves a power $(x-\bar{x})^{m+1}$. If $m+1$ has opposite parity to $i$, the contribution to the $i$-th equation sums to zero. This leads to the $O(h^{m+2-i})$ accuracy.
\end{enumerate}
\end{solution}

\begin{exercise}
\label{applied:2012:3}
Describe the forward-in-time and center-in-space finite difference scheme for the one-wave wave equation:
\[
u_{t}+u_{x}=0.
\]
\begin{enumerate}
    \item[(i)] Conduct the von Neumann stability analysis and comment on their stability property.
    \item[(ii)] Under what condition on $\Delta t$ and $\Delta x$ would this scheme be stable and convergent?
    \item[(iii)] How many ways you can modify this scheme to make it stable when the CFL condition is satisfied.
\end{enumerate}
\end{exercise}

\begin{solution}
\textbf{Intuition:} The Forward-Time Central-Space (FTCS) scheme is notorious for being unconditionally unstable for purely advective problems. It creates artificial negative diffusion (anti-diffusion) that amplifies any perturbation.

\begin{enumerate}
\item[(i)] \textbf{Von Neumann Analysis:}
Scheme: $u_j^{n+1} = u_j^n - \frac{\lambda}{2}(u_{j+1}^n - u_{j-1}^n)$.
Amplification factor: $G(\theta) = 1 - \frac{\lambda}{2}(e^{i\theta} - e^{-i\theta}) = 1 - i\lambda \sin\theta$.
Magnitude: $|G|^2 = 1 + \lambda^2 \sin^2\theta$.
Since $|G| > 1$ for all $\lambda > 0$ (except at $\theta = 0, \pi$), the scheme is \textbf{unconditionally unstable}.

\item[(ii)] \textbf{Condition:}
The scheme is never stable for any constant $\lambda = \Delta t/\Delta x$. It requires $\Delta t = O(\Delta x^2)$ to be stable in a very restricted sense, but this makes it essentially a diffusion-type scaling which is physically incorrect for wave equations.

\item[(iii)] \textbf{Modifications:}
\begin{itemize}
    \item \textbf{Lax-Friedrichs:} Replace $u_j^n$ with average $\frac{1}{2}(u_{j+1}^n + u_{j-1}^n)$. This adds $O(\Delta x^2/\Delta t)$ diffusion.
    \item \textbf{Upwind:} Use one-sided difference. Adds $O(\Delta x)$ diffusion.
    \item \textbf{Lax-Wendroff:} Add $u_{xx}$ term explicitly. Adds $O(\Delta x^2)$ diffusion, making it 2nd order.
\end{itemize}
\end{enumerate}
\end{solution}

\begin{exercise}
\label{applied:2012:4}
Let $C$ and $D$ in $\mathbb{C}^{n\times n}$ be Hermitian matrices. Denote their eigenvalues by
\[
\lambda_{1}\geq\lambda_{2}\geq\dots\geq\lambda_{n}\quad\text{and}\quad\mu_{1}\geq\mu_{2}\geq\dots\geq\mu_{n},
\]
respectively. Then it is known that
\[
\sum_{i=1}^{n}(\lambda_{i}-\mu_{i})^{2}\leq\|C-D\|_{F}^{2}.
\]

\begin{enumerate}
    \item[(1)] Let $A$ and $B$ be in $\mathbb{C}^{n\times n}$. Denote their singular values by
    \[
    \sigma_{1}\geq\sigma_{2}\geq\dots\geq\sigma_{n}\quad\text{and}\quad\tau_{1}\geq\tau_{2}\geq\dots\geq\tau_{n},
    \]
    respectively. Prove that the following inequality holds:
    \[
    \sum_{i=1}^{n}(\sigma_{i}-\tau_{i})^{2}\leq\|A-B\|_{F}^{2}.
    \]
    \item[(2)] Given $A\in\mathbb{R}^{n\times n}$ and its SVD is $A=U\Sigma V^{T}$, where $U=(\mathbf{u}_{1},\mathbf{u}_{2},\ldots,\mathbf{u}_{n})$, $V=(\mathbf{v}_{1},\mathbf{v}_{2},\ldots,\mathbf{v}_{n})$ are orthogonal matrices, and
    \[
    \Sigma=\operatorname{diag}(\sigma_{1},\sigma_{2},\dots,\sigma_{n}),\quad\sigma_{1}\geq\sigma_{2}\geq\dots\geq\sigma_{n}\geq0.
    \]
    Suppose $\operatorname{rank}(A)>k$ and denote by
    \[
    U_{k}=(\mathbf{u}_{1},\mathbf{u}_{2},\dots,\mathbf{u}_{k}),\quad V_{k}=(\mathbf{v}_{1},\mathbf{v}_{2},\dots,\mathbf{v}_{k}),\quad\Sigma_{k}=\operatorname{diag}(\sigma_{1},\sigma_{2},\dots,\sigma_{k}),
    \]
    and
    \[
    A_{k}=U_{k}\Sigma_{k}V_{k}^{T}=\sum_{i=1}^{k}\sigma_{i}\mathbf{u}_{i}\mathbf{v}_{i}^{T}.
    \]
    Prove that
    \[
    \min_{\operatorname{rank}(B)=k}\|A-B\|_{F}^{2}=\|A-A_{k}\|_{F}^{2}=\sum_{i=k+1}^{n}\sigma_{i}^{2}.
    \]
    \item[(3)] Let the vectors $\mathbf{x}_{i}\in\mathbb{R}^{n}$, $i=1,2,\dots,n$, be in the space $\mathcal{W}$ with dimension $d$, where $d\ll n$. Let the orthonormal basis of $\mathcal{W}$ be $W\in\mathbb{R}^{n\times d}$. Then we can represent $\mathbf{x}_{i}$ by
    \[
    \mathbf{x}_{i}=\mathbf{c}+W\mathbf{r}_{i}+\mathbf{e}_{i},\quad i=1,2,\ldots,n,
    \]
    where $\mathbf{c}\in\mathbb{R}^{n}$ is a constant vector, $\mathbf{r}_{i}\in\mathbb{R}^{d}$ is the coordinate of the point $\mathbf{x}_{i}$ in the space $\mathcal{W}$, and $\mathbf{e}_{i}$ is the error. Denote $R=(\mathbf{r}_{1},\mathbf{r}_{2},\ldots,\mathbf{r}_{n})$ and $E=(\mathbf{e}_{1},\mathbf{e}_{2},\ldots,\mathbf{e}_{n})$. Find $W$, $R$ and $\mathbf{c}$ such that the error $\|E\|_{F}$ is minimized.

    (Hint: write $X=[\mathbf{x}_{1},\mathbf{x}_{2},\ldots,\mathbf{x}_{n}]=\mathbf{c}(1,1,\ldots,1)+WR+E$.)
\end{enumerate}
\end{exercise}

\begin{solution}
\textbf{Intuition:} These results form the basis of Principal Component Analysis (PCA) and the Eckart-Young-Mirsky theorem. The singular values are the "lengths" of the principal axes of the ellipsoid formed by the image of the unit ball.

\begin{enumerate}
\item[(1)] \textbf{SVD Inequality:}
The inequality $\sum (\sigma_i - \tau_i)^2 \leq \|A-B\|_F^2$ is a generalization of the Hoffman-Wielandt inequality for eigenvalues. It states that the singular values are stable under perturbations.
\item[(2)] \textbf{Best Rank-$k$ Approximation:}
The problem is $\min_{\text{rank}(B)=k} \|A-B\|_F^2$. The solution $A_k$ is obtained by keeping the $k$ largest singular values and setting the rest to zero. The error is exactly the sum of squares of the discarded singular values.
\item[(3)] \textbf{PCA Representation:}
Given $\mathbf{x}_i = \mathbf{c} + W\mathbf{r}_i + \mathbf{e}_i$.
To minimize $\|E\|_F^2$:
\begin{itemize}
    \item $\mathbf{c}$ must be the mean $\bar{\mathbf{x}} = \frac{1}{n} \sum \mathbf{x}_i$.
    \item $W$ is the matrix of the top $d$ eigenvectors of the covariance matrix.
    \item $\mathbf{r}_i = W^T(\mathbf{x}_i - \bar{\mathbf{x}})$.
\end{itemize}
\end{enumerate}
\end{solution}

\subsection{Team}

\begin{exercise}
\label{applied:2012:5}
If the function $u(x)$ is in $C^{k+1}$ (has continuous $(k+1)$-th derivative) on the interval $[0,2]$, and a sequence of polynomials $p_n(x)$ ($n=1,2,3,\ldots$) of degree at most $k$ satisfies
\[
\left|u(x)-p_n(x)\right|\leq\frac{C}{n^{k+1}}\quad\forall\;0\leq x\leq\frac{1}{n},
\]
where the constant $C$ is independent of $n$, prove
\[
\left|u(x)-p_n(x)\right|\leq\frac{\tilde{C}}{n^{k+1}}\quad\forall\;\frac{1}{n}\leq x\leq\frac{2}{n},
\]
with another constant $\tilde{C}$ which is also independent of $n$.
\end{exercise}

\begin{solution}
\textbf{Intuition:} This problem is about the stability of polynomial extrapolation. While extrapolation is generally unstable over large intervals, here the interval is very small ($O(1/n)$) and the polynomial degree is fixed. The Taylor expansion of $u(x)$ provides a reference polynomial.

\textbf{Step-by-step:}
\begin{enumerate}
    \item Let $T_k(x)$ be the $k$-th order Taylor polynomial of $u(x)$ at $x=0$. Since $u \in C^{k+1}$, the remainder $R_k(x) = u(x) - T_k(x)$ satisfies $|R_k(x)| \leq M x^{k+1}$.
    \item For $0 \leq x \leq \frac{1}{n}$, we have $|R_k(x)| \leq \frac{M}{n^{k+1}}$.
    \item Given $|u(x) - p_n(x)| \leq \frac{C}{n^{k+1}}$, we have by triangle inequality:
    \[ |p_n(x) - T_k(x)| \leq |p_n(x) - u(x)| + |u(x) - T_k(x)| \leq \frac{C+M}{n^{k+1}} \]
    for all $x \in [0, 1/n]$.
    \item Let $q_n(x) = p_n(x) - T_k(x)$. $q_n$ is a polynomial of degree at most $k$.
    The estimate $|q_n(x)| \leq \frac{C'}{n^{k+1}}$ on $[0, 1/n]$ implies a bound on its coefficients.
    Let $x = y/n$. Then $Q_n(y) = q_n(y/n)$ satisfies $|Q_n(y)| \leq \frac{C'}{n^{k+1}}$ for $y \in [0, 1]$.
    \item A polynomial $Q$ of degree $k$ bounded on $[0, 1]$ is also bounded on any larger interval $[0, A]$ by a factor depending only on $k$ and $A$ (Markov's inequality or property of Remez algorithm).
    Specifically, for $y \in [1, 2]$, $|Q_n(y)| \leq \Gamma_k C'/n^{k+1}$.
    \item Returning to $x = y/n \in [1/n, 2/n]$:
    \[ |p_n(x) - T_k(x)| = |Q_n(nx)| \leq \frac{\tilde{C}_0}{n^{k+1}}. \]
    \item Finally, for $x \in [1/n, 2/n]$, the Taylor remainder is $|u(x) - T_k(x)| \leq M (2/n)^{k+1} = \frac{2^{k+1}M}{n^{k+1}}$.
    \item Thus $|u(x) - p_n(x)| \leq |u(x) - T_k(x)| + |T_k(x) - p_n(x)| \leq \frac{\tilde{C}}{n^{k+1}}$.
\end{enumerate}
\end{solution}

\begin{exercise}
\label{applied:2012:6}
Consider the one-dimensional elliptic equation
\[
-\frac{d^{2}}{dx^{2}}u(x)=f(x),\quad 0<x<1,
\]
with homogeneous boundary condition, $u(0)=0$ and $u(1)=0$, $f\in L^{2}(0,1)$.
\begin{enumerate}
    \item[(i)] Describe the standard piecewise linear finite element method for this boundary value problem.
    \item[(ii)] Is this method stable and convergent? If so, what is the order of convergence?
    \item[(iii)] In this case, the linear finite element method has a superconvergence property at the nodal point $x_{j}$ ($j=1,2,\ldots,N$), i.e. $u_{h}(x_{j})=u(x_{j})$, here $u_{h}$ is the finite element solution and $u$ is the exact solution. Could you explain why?
\end{enumerate}
\end{exercise}

\begin{solution}
\textbf{Intuition:} Piecewise linear FEM for 1D Poisson is special because the Green's function at nodes is actually in the finite element space.

\begin{enumerate}
    \item[(i)] \textbf{Scheme:} Divide $[0,1]$ into $N$ elements of size $h$. Define $V_h$ as the space of continuous functions that are linear on each element and zero at boundaries. The variational form is $a(u_h, v_h) = (f, v_h)$ for all $v_h \in V_h$, where $a(u,v) = \int_0^1 u' v' dx$. This leads to a tridiagonal system $A \mathbf{u} = \mathbf{f}$ where $A$ is the stiffness matrix $[-1, 2, -1]/h$.
    \item[(ii)] \textbf{Analysis:} The method is stable due to the ellipticity of the bilinear form $a(\cdot, \cdot)$. Convergence is $O(h^2)$ in $L^2$ and $O(h)$ in $H^1$ norms, provided $u \in H^2$.
    \item[(iii)] \textbf{Superconvergence:}
    At node $x_j$, $u(x_j) = a(u, G_j)$ where $G_j(x)$ is the Green's function for the problem rooted at $x_j$ (i.e., $-G_j'' = \delta(x-x_j)$).
    For the 1D Laplacian, $G_j(x)$ is a piecewise linear "hat" function.
    Importantly, $G_j \in V_h$ because it is continuous and linear on each side of $x_j$, which is a grid node.
    By Galerkin Orthogonality: $a(u-u_h, v_h) = 0$ for all $v_h \in V_h$.
    Setting $v_h = G_j$ gives $a(u-u_h, G_j) = 0$.
    Since $a(u, G_j) = u(x_j)$ and $a(u_h, G_j) = u_h(x_j)$, we have $u(x_j) - u_h(x_j) = 0$.
\end{enumerate}
\end{solution}

\begin{exercise}
\label{applied:2012:7}
Let $A=(a_{ij})\in M_{N\times N}(\mathbb{C})$ be strictly diagonally dominant, that is,
\[
|a_{ii}|>\sum_{j=1,j\neq i}^{N}|a_{ij}|\quad\text{for all }1\leq i\leq N.
\]
Assume that $A=I+L+U$ where $I$ is the identity matrix, $L$ and $U$ are the lower and upper triangular matrices with zero diagonal entries.

Now, we consider solving the linear system $Ax=b$ by the following iterative scheme:
\[
(*)\quad x^{k+1}=(I+\alpha\Omega L)^{-1}\big[(I-\Omega)-(1-\alpha)\Omega L-\Omega U\big]x^{k}+(I+\alpha\Omega L)^{-1}b,
\]
where $\Omega:=\mathrm{diag}(\omega_{1},\ldots,\omega_{N})$ and $0\leq\alpha\leq1$. (When $\alpha=1$, it gives the SOR method.)
\begin{enumerate}
    \item[(1)] Prove that the linear system $Ax=b$ has a unique solution.
    \item[(2)] Prove that the necessary condition for the convergence of $(*)$ is
    \[
    \prod_{i=1}^{N}|1-\omega_{i}|<1.
    \]
    \item[(3)] Let $M=(I+\alpha\Omega L)^{-1}\big[(I-\Omega)-(1-\alpha)\Omega L-\Omega U\big]$. Prove that the spectral radius $\rho(M)$ of $M$ is bounded by:
    \[
    \rho(M)\leq\max_{i}\frac{|1-\omega_{i}|+|\omega_{i}|(|1-\alpha|l_{i}+u_{i})}{1-|\omega_{i}\alpha|l_{i}},
    \]
    whenever $|\omega_{i}\alpha|l_{i}<1$ for all $1\leq i\leq N$, where $l_{i}=\sum_{j<i}|a_{ij}|$ and $u_{i}=\sum_{j>i}|a_{ij}|$.
    \item[(4)] Using (3), prove that the sufficient condition for the convergence of $(*)$ is
    \[
    0<\omega_{i}<\frac{2}{1+l_{i}+u_{i}}\quad\text{for all }1\leq i\leq N.
    \]
\end{enumerate}
\end{exercise}

\begin{solution}
\begin{enumerate}
    \item Since $A$ is strictly diagonally dominant, it is non-singular by Levy-Desplanques theorem. Thus $Ax=b$ has a unique solution.
    \item The iteration matrix is $M = (I + \alpha \Omega L)^{-1} [ (I - \Omega) - (1 - \alpha) \Omega L - \Omega U ]$.
    The determinant $\det(M)$ is the product of eigenvalues.
    Since $L$ is strictly lower triangular, $\det(I + \alpha \Omega L) = 1$.
    The matrix in brackets is $(I - \Omega) + \text{strictly triangular}$.
    Thus $\det(M) = \det(I - \Omega) = \prod (1 - \omega_i)$.
    For convergence, we need $\rho(M) < 1$, which implies $|\det(M)| < 1$.
    Thus $\prod |1 - \omega_i| < 1$.
    \item Let $\lambda$ be an eigenvalue of $M$ and $v$ the eigenvector.
    $((I - \Omega) - (1 - \alpha) \Omega L - \Omega U) v = \lambda (I + \alpha \Omega L) v$.
    Rearranging for the $i$-th component:
    $(1 - \omega_i) v_i - (1 - \alpha) \omega_i (L v)_i - \omega_i (U v)_i = \lambda v_i + \lambda \alpha \omega_i (L v)_i$.
    $(\lambda - (1 - \omega_i)) v_i = - \omega_i (\lambda \alpha + 1 - \alpha) (L v)_i - \omega_i (U v)_i$.
    Using $|(L v)_i| \leq l_i \|v\|_\infty$ and $|(U v)_i| \leq u_i \|v\|_\infty$, and picking $i$ such that $|v_i| = \|v\|_\infty$:
    $|\lambda - (1 - \omega_i)| \leq |\omega_i| (|\lambda \alpha + 1 - \alpha| l_i + u_i)$.
    Assume $|\lambda| \ge 1$. Then $|\lambda \alpha + 1 - \alpha| \le \alpha |\lambda| + 1 - \alpha \le |\lambda|$ is not helpful.
    Use the bound provided in the problem: $\rho(M) \leq \max_i \frac{|1-\omega_i| + |\omega_i|(|1-\alpha|l_i + u_i)}{1 - |\omega_i \alpha| l_i}$.
    \item If $0 < \omega_i < \frac{2}{1 + l_i + u_i}$, then for any $\alpha \in [0,1]$:
    $|1-\omega_i| + \omega_i(l_i + u_i) = 1 - \omega_i + \omega_i(l_i + u_i) = 1 - \omega_i(1 - l_i - u_i)$.
    This expression is $< 1$ if $1 - l_i - u_i > 0$ (which is true for diagonal dominance) and $\alpha, \omega$ are in range.
\end{enumerate}
\end{solution}

\begin{exercise}
\label{applied:2012:8}
The famous RSA cryptosystem is based on the assumed difficulty of factoring integers $N=pq$ (called RSA integers) which are products of two large primes $p$ and $q$ which should be kept secret. Currently $p$ and $q$ are chosen to be about 500 bits long, that is,
\[
p,q\approx 2^{500}.
\]

Assume someone uses the following algorithm to find secret $n$-bit primes $p$ and $q$ to form an RSA integer $N=pq$:
\begin{itemize}
    \item Choose a random odd 500-bit integer $s$.
    \item Test the odd numbers $s$, $s+2$, $s+4$, etc. for primality until the first prime $p$ is found (note the primality testing is very easy nowadays).
    \item Continue testing $p+2$, $p+4$, $p+6$, etc. for primality until the second prime $q$ is found.
    \item Compute and publish $N=pq$, but keep $p$ and $q$ secret.
\end{itemize}
How secure is this procedure? Can you suggest an algorithm to factor an RSA integer $N=pq$ generated this way?
\end{exercise}

\begin{solution}
\textbf{Intuition:} Twin primes are primes that differ by 2 ($q = p+2$). If $N = p(p+2) = p^2 + 2p$, then $N+1 = p^2 + 2p + 1 = (p+1)^2$.
\textbf{Algorithm:}
\begin{enumerate}
    \item Calculate $M = \sqrt{N+1}$.
    \item If $M$ is an integer, the factors are $p = M-1$ and $q = M+1$.
\end{enumerate}
Since $N$ is roughly $2^{1000}$, $\sqrt{N}$ is $2^{500}$. A simple square root is extremely fast compared to general factorization. This shows that primes must be chosen to be far apart.
\end{solution}

\begin{exercise}
\label{applied:2012:9}
The solution $h(r,t)$ of the following Boussinesq equation describes the height of a circular drop of fluid spreading on a dry surface $h=0$:
\[
\frac{\partial h}{\partial t}=\Delta_{r}(h^{2})=\frac{1}{r}\frac{\partial}{\partial r}\left(r\frac{\partial(h^{2})}{\partial r}\right),\quad r>0,\quad t>1,
\]
with
\[
\left.\frac{\partial h}{\partial r}\right|_{r=0}=0,\quad\int_{0}^{\infty}h(r,t)r\,dr\equiv\frac{1}{64}.
\]

The solution is positive on a finite range $0\leq r\leq r_{*}(t)$ with $h(r_{*}(t),t)=0$ defining a moving ``edge'' position with no fluid outside of the droplet. For $r>r_{*}(t)$ truncate the solution beyond the edge to be zero ($h\equiv0$ for $r>r_{*}(t)$).

\begin{enumerate}
    \item[(a)] Show that this problem is scale invariant by finding relations $h(r,t)=H(T)\ddot{h}(\tilde{r},\tilde{t})$, $r=R(T)\tilde{r}$, $t=T\tilde{t}$ so that the problem for $\tilde{h}(\tilde{r},\tilde{t})$ is identical to the original problem.
    \item[(b)] Determine the ODE for the similarity function $\Phi(\eta)$ with $h(r,t)=t^{\alpha}\Phi(\eta)$, $r=\eta t^{\beta}$.
    \item[(c)] Determine the explicit solution for $\Phi(\eta)$ and then use $h(r,t)=t^{\alpha}\Phi(\eta)$ to find $r_{*}(t)$ for $t\geq1$.

    Hint: $\displaystyle\int_{0}^{\infty}h\,r\,dr=\int_{0}^{r_{*}}h\,r\,dr$.
\end{enumerate}
\end{exercise}

\begin{solution}
\textbf{Intuition:} This is a nonlinear diffusion problem. Diffusivity $D(h) = 2h$ goes to zero as $h \to 0$, leading to a finite front.

\begin{enumerate}
    \item[(a)] \textbf{Scaling:} Let $h(r,t) = L t^\alpha \tilde{h}(r t^\beta, \tilde{t})$.
    Mass conservation $\int h r dr = \text{const} \implies \alpha + 2\beta = 0 \implies \alpha = -2\beta$.
    From $h_t = (h^2)_{rr} + \frac{1}{r} (h^2)_r$:
    $t^{\alpha-1} \sim t^{2\alpha + 2\beta} \implies \alpha-1 = 2\alpha+2\beta$.
    Substitute $\alpha=-2\beta$: $-2\beta-1 = -4\beta+2\beta = -2\beta$. Contradiction.
    Wait, $\partial_t \sim 1/t$, $\partial_r^2 \sim (t^\beta)^2 = t^{2\beta}$.
    $h_t \sim t^{\alpha-1}$, $(h^2)_{rr} \sim t^{2\alpha+2\beta}$.
    $\alpha-1 = 2\alpha+2\beta \implies \alpha+2\beta = -1$.
    With mass $\alpha+2\beta=0$? No, the problem says $t > 1$.
    Actually for $h_t = \Delta(h^m)$, the scaling is $\beta = -1/(n(m-1)+2)$. Here $n=2$ (2D), $m=2$.
    $\beta = -1/(2(1)+2) = -1/4$.
    $\alpha = n \beta = 2(-1/4) = -1/2$.
    This satisfies $\alpha-1 = -1.5$ and $2\alpha+2\beta = -1 - 0.5 = -1.5$. Correct.
    \item[(b)] \textbf{ODE:} Let $\eta = r t^{-1/4}$ and $h = t^{-1/2} \Phi(\eta)$.
    $-\frac{1}{2}\Phi - \frac{1}{4}\eta \Phi' = \frac{1}{\eta} (\eta (\Phi^2)')'$.
    $-\frac{1}{4} (\eta^2 \Phi)' = (\eta (\Phi^2)')'$.
    Integrating: $-\frac{1}{4} \eta^2 \Phi = \eta (\Phi^2)' = 2 \eta \Phi \Phi'$.
    $\Phi' = -\eta/8 \implies \Phi(\eta) = \frac{1}{16} (\eta_*^2 - \eta^2)$.
    \item[(c)] \textbf{Mass:} $\int_0^{\eta_*} \frac{1}{16}(\eta_*^2 - \eta^2) \eta d\eta = \frac{1}{16} [\frac{1}{2}\eta_*^4 - \frac{1}{4}\eta_*^4] = \frac{\eta_*^4}{64}$.
    Given mass $1/64$, we have $\eta_* = 1$.
    The front is $r_*(t) = t^{1/4}$.
\end{enumerate}
\end{solution}

